\documentclass[11pt]{ctexart}
\usepackage[a4paper,margin=28mm,headheight=14pt]{geometry}
\usepackage{amsmath,amssymb,amsthm,mathtools,enumitem,microtype,hyperref,fancyhdr}
\hypersetup{colorlinks=true,linkcolor=blue,citecolor=blue,urlcolor=blue}
\setlength{\parskip}{0.35em}
\setlength{\parindent}{2em}
\emergencystretch=3em
\sloppy
\setlist{nosep,leftmargin=2em}
\pagestyle{fancy}\fancyhf{}\fancyhead[L]{Carlet 猜想的证明候选稿}\fancyhead[R]{王泉智、ChatGPT}\fancyfoot[C]{\thepage}
\title{Carlet 猜想的一个完整证明候选\\[0.4em]\large Reciprocal 紧化、二次完全交与显式 Betti 界}
\author{王泉智\thanks{通讯邮箱：\href{mailto:mrhwrn120504@gmail.com}{mrhwrn120504@gmail.com}} \and ChatGPT（AI 辅助写作与编辑）}
\date{2026 年 11 月\quad（基于 2026 年 8 月稿件修订）}
\begin{document}\maketitle
\begin{center}\fbox{\begin{minipage}{0.92\linewidth}\small\textbf{稿件状态。} 本文整理并补充一份完整证明候选稿。三份有限域证书已有独立精确验证；第 4--10 节的统一 reciprocal 几何论证尚未完成独立专家逐行复核。本文不把该候选证明表述为已经通过 Lean 或独立审稿认证的定理。\end{minipage}}\end{center}
\begin{abstract}
本文整理一条关于二元有限域乘法逆函数零和子空间的完整证明候选路线。对奇数 $n$，目标是对每个 $3\le k\le n-3$ 构造 $k$ 维 $\mathbb F_2$-子空间 $E\subset\mathbb F_{2^n}$，使 $\sum_{x\in E\setminus\{0\}}x^{-1}=0$。证明路线先用子空间多项式把零和条件化为 $X^2$ 系数消失，再将其表示为 Pink--Schieder reciprocal 紧化中的 Cartier 截面；随后给出满秩层的显式二次模型、分层上同调估计和有限尾部证书。附录的 $(n,k)=(29,13),(29,14),(31,14)$ 证书已由独立精确计算复核。统一几何部分仍待独立逐行审查，因此本文准确地称为完整证明候选稿，而非已认证证明。
\end{abstract}
\noindent\textbf{关键词：} Carlet 猜想；有限域；乘法逆函数；子空间多项式；reciprocal plane；交叉上同调；Betti 数。\medskip
\tableofcontents\clearpage
\section{引言与候选主定理}

\subsection{Carlet 猜想}

令 \ensuremath{\mathbb F_{2^n}} 为含 \ensuremath{2^{n}} 个元素的二元有限域，并令 f\_inv(0)=0、f\_inv(x)=\ensuremath{x^{-1}}（x\ensuremath{\ne}0）。Carlet 引入 k-th order sum-free 的概念，并证明对乘法逆函数而言，相关维数上的失败可以用线性子空间上的逆元零和来刻画。本文使用下列等价形式作为工作定义。

\begin{quote}\textbf{猜想 1.1（Carlet）. 若 n 为奇数，则对每个整数 3\ensuremath{\le}k\ensuremath{\le}n−3，存在 k 维 \ensuremath{F_2}-线性子空间 E\ensuremath{\subset}\ensuremath{\mathbb F_{2^n}}，满足 \ensuremath{\Sigma}\_\{0\ensuremath{\ne}x\ensuremath{\in}E\}\ensuremath{x^{-1}}=0。}\end{quote}

记 \ensuremath{K_{n}} 为满足上述存在性的维数集合。Carlet–Hou 及后续工作给出互补维数对称 k\ensuremath{\in}\ensuremath{K_{n}} \ensuremath{\Longleftrightarrow} n−k\ensuremath{\in}\ensuremath{K_{n}}，因此奇数 n 时只需处理 3\ensuremath{\le}k\ensuremath{\le}(n−1)/2。

\subsection{本文结果的状态}

\begin{quote}\textbf{定理 1.2（已验证有限结果）. Carlet 猜想对 n=29 和 n=31 成立。更具体地，本文附录 A 给出 (29,13)、(29,14)、(31,14) 的显式零和子空间；结合公开结果与互补维数对称即可覆盖全部允许维数。}\end{quote}

\begin{quote}\textbf{定理 1.2 的有限域证书已经由两种独立方式精确验证：完整子空间多项式的 \ensuremath{X^2} 系数为零，以及对全部非零子空间元素逐个求逆后直接求和为零。附带脚本 verify\_carlet\_paper.py 可重新计算这些事实。}\end{quote}

\begin{quote}\textbf{候选主定理 1.3. 对每个奇数 n 和每个 3\ensuremath{\le}k\ensuremath{\le}n−3，存在 k 维 \ensuremath{F_2}-子空间 E\ensuremath{\subset}\ensuremath{\mathbb F_{2^n}} 使 \ensuremath{\Sigma}\_\{0\ensuremath{\ne}x\ensuremath{\in}E\}\ensuremath{x^{-1}}=0。  [待独立复核]}\end{quote}

本文第 4–10 节给出候选主定理 1.3 的完整推导。为避免研究状态与逻辑内容混淆，凡依赖本文新建 reciprocal 几何桥梁的结论均应理解为“候选证明中的定理”；外部已发表定理和有限证书则按通常定理使用。

\subsection{证明策略}

\begin{itemize}

\item 把零和条件写成 reciprocal coordinates 下的线性方程 \ensuremath{\ell}\_k=0，并在 Pink–Schieder compactification \ensuremath{Q_{k}} 中定义 Cartier divisor \ensuremath{D_{k}}。

\item 证明满秩开层 \ensuremath{X_{r}} 与一个显式二次 affine scheme \ensuremath{Y_{r}} 在 scheme 层面同构；\ensuremath{Y_{r}} 是光滑的 set-theoretic complete intersection。

\item 用 Pink–Schieder 的 rank-stratification 和局部乘积、EPW 的 modular-matroid reciprocal-plane intersection cohomology，证明 \ensuremath{Q_{k}} 与 \ensuremath{D_{k}} rationally smooth。

\item 用 Pink–Schieder strange morphisms 得到 \ensuremath{Q_{k}} 的 Frobenius-equivariant \ensuremath{\ell}-adic 上同调等同射影空间；再用 Wan–Zhang weak Lefschetz 与 Verdier 对偶把 \ensuremath{D_{k}} 的未知上同调压缩到唯一 middle degree。

\item 用 Wan–Zhang 对 affine complete intersection 的 Betti 上界控制 middle error；用一个次数 \ensuremath{2^{r}}−2 的清分母多项式和 Schwartz–Zippel 控制低秩边界。

\item 证明大参数主项严格压过 middle error 与 boundary；最后用公开小维结果和 n=29、31 证书闭合有限尾部。

\end{itemize}

\section{预备知识：有限域、子空间多项式与互补对称}

\paragraph{补充说明（任意基环上的定义）.} 对任意交换 $\mathbb F_2$-代数 $A$ 及元素 $x_1,\ldots,x_j\in A$，先定义形式多项式 $L_j(T)=\prod_{\epsilon\in\mathbb F_2^j}(T+\sum_i\epsilon_i x_i)\in A[T]$。这一乘积在尚未证明坐标线性无关时仍有定义；只有在单位条件推出所有非零线性组合均为单位之后，才把坐标解释为 Drinfeld 开集上的点。

\subsection{子空间多项式}

\begin{quote}\textbf{定义 2.1. 对有限 \ensuremath{F_2}-子空间 E，定义 \ensuremath{L_{E}}(X)=\ensuremath{\prod}\_\{u\ensuremath{\in}E\}(X+u)。若 dim\_\{\ensuremath{F_2}\}E=r，则 \ensuremath{L_{E}} 是 monic 线性化多项式，形如 \ensuremath{L_{E}}(X)=\ensuremath{c_0}X+\ensuremath{c_1}\ensuremath{X^2}+\ensuremath{\cdots}+c\_rX\textasciicircum\{\}\{\ensuremath{2^{r}}\}，其中 \ensuremath{c_0}\ensuremath{\ne}0、\ensuremath{c_{r}}=1。}\end{quote}

\begin{quote}\textbf{引理 2.2（逆元和判据）. 设 E\ensuremath{\ne}\{0\}。则 \ensuremath{\Sigma}\_\{0\ensuremath{\ne}u\ensuremath{\in}E\}\ensuremath{u^{-1}}=\ensuremath{c_1}/\ensuremath{c_0}。特别地，E 为零和子空间当且仅当 \ensuremath{c_1}=0。}\end{quote}

\begin{proof}[证明] 

令 \ensuremath{Q_{E}}(X)=\ensuremath{L_{E}}(X)/X=\ensuremath{\prod}\_\{0\ensuremath{\ne}u\ensuremath{\in}E\}(X+u)。由于 \ensuremath{Q_{E}}(0)=\ensuremath{\prod}\_\{0\ensuremath{\ne}u\ensuremath{\in}E\}u=\ensuremath{c_0}\ensuremath{\ne}0，对 \ensuremath{Q_{E}} 作对数导数可得 \ensuremath{Q_{E}}\ensuremath{\prime}(0)/\ensuremath{Q_{E}}(0)=\ensuremath{\Sigma}\_\{0\ensuremath{\ne}u\ensuremath{\in}E\}\ensuremath{u^{-1}}。另一方面，由线性化形式可写 \ensuremath{Q_{E}}(X)=\ensuremath{c_0}+\ensuremath{c_1}X+\ensuremath{c_2}\ensuremath{X^3}+\ensuremath{\cdots}，故 \ensuremath{Q_{E}}\ensuremath{\prime}(0)=\ensuremath{c_1}。代入即得。

\end{proof}

\begin{quote}\textbf{引理 2.3（扩维递推）. 若 v\ensuremath{\notin}E 且 E\ensuremath{^+}=E\ensuremath{\oplus}\ensuremath{\langle}v\ensuremath{\rangle}，则 \ensuremath{L_{E+}}(X)=\ensuremath{L_{E}}(X)²+\ensuremath{L_{E}}(v)\ensuremath{L_{E}}(X)。}\end{quote}

\begin{proof}[证明] 

因为 E\ensuremath{^+}=E \ensuremath{\sqcup} (E+v)，故 \ensuremath{L_{E+}}(X)=\ensuremath{L_{E}}(X)\ensuremath{L_{E}}(X+v)。线性化多项式满足 \ensuremath{L_{E}}(X+v)=\ensuremath{L_{E}}(X)+\ensuremath{L_{E}}(v)，于是结论成立。

\end{proof}

\begin{quote}\textbf{推论 2.4（最后一步线性求解）. 若 \ensuremath{L_{E}}(X)=\ensuremath{c_0}X+\ensuremath{c_1}\ensuremath{X^2}+\ensuremath{\cdots} 且 \ensuremath{c_1}\ensuremath{\ne}0，则 E\ensuremath{\oplus}\ensuremath{\langle}v\ensuremath{\rangle} 为零和空间当且仅当 \ensuremath{L_{E}}(v)=\ensuremath{c_0}²/\ensuremath{c_1}。}\end{quote}

\begin{proof}[证明] 

由引理 2.3，扩维后 \ensuremath{X^2} 系数为 \ensuremath{c_1}\ensuremath{^+}=\ensuremath{c_0}²+\ensuremath{L_{E}}(v)\ensuremath{c_1}。结合引理 2.2 即得。由于 \ensuremath{L_{E}} 是 \ensuremath{F_2}-线性映射，该条件是一个普通二元线性方程。

\end{proof}

\subsection{互补维数}

本文将 Carlet–Hou/Hou–Zhao 的互补维数定理作为外部输入：k\ensuremath{\in}\ensuremath{K_{n}} 当且仅当 n−k\ensuremath{\in}\ensuremath{K_{n}}。因而对奇数 n，只需证明低半区 3\ensuremath{\le}k\ensuremath{\le}(n−1)/2。该区间等价于 n\ensuremath{\ge}2k+1，这正是后文点数估计的自然参数条件。

\section{Reciprocal compactification 与 Carlet divisor}

\subsection{Pink–Schieder 的 reciprocal ring}

令 V=\ensuremath{F_2}\textasciicircum\{\}k。Pink–Schieder 在对称代数 \ensuremath{S_{V}}=Sym(V) 的分式域 \ensuremath{K_{V}} 中定义 reciprocal ring

\ensuremath{R_{V}} = \ensuremath{F_2}[ 1/v : 0\ensuremath{\ne}v\ensuremath{\in}V ] \ensuremath{\subset} \ensuremath{K_{V}}.    (3.1)

按 Pink–Schieder 的原始约定，1/v 的次数为 −1；本文反转 grading，使 reciprocal generators 的次数为 +1，并记

\ensuremath{Q_{V}} = Proj \ensuremath{R_{V}}.    (3.2)

Pink–Schieder [PS, Thm. 1.7, 1.11] 证明 \ensuremath{R_{V}} 为 Cohen–Macaulay normal integral domain，\ensuremath{Q_{V}} 为 integral、Cohen–Macaulay、projectively normal。开层 \ensuremath{\Omega}\_V 同时嵌入 P(V) 与 \ensuremath{Q_{V}}，并等同于 \ensuremath{P^{k−1}} 删除所有 \ensuremath{F_2}-rational hyperplanes 后的 Drinfeld period domain。

\subsection{Carlet 线性截面}

对每个 0\ensuremath{\ne}v\ensuremath{\in}V，记 \ensuremath{z_{v}}=1/v 为 reciprocal homogeneous coordinate。定义

\ensuremath{\ell}\_k = \ensuremath{\Sigma}\_\{0\ensuremath{\ne}v\ensuremath{\in}V\} \ensuremath{z_{v}},    \ensuremath{D_{k}} = V(\ensuremath{\ell}\_k) \ensuremath{\subset} \ensuremath{Q_{V}},    \ensuremath{X_{k}} = \ensuremath{D_{k}} \ensuremath{\cap} \ensuremath{\Omega}\_V.    (3.3)

\begin{quote}\textbf{引理 3.1. \ensuremath{\ell}\_k 不是 \ensuremath{R_{V}} 中的零元。因此 \ensuremath{D_{k}} 是 \ensuremath{Q_{V}} 上由 O(1) 的非零 section 定义的 effective ample Cartier divisor，且 dim \ensuremath{D_{k}}=k−2。}\end{quote}

\begin{proof}[证明] 

只需证明 \ensuremath{\ell}\_k 作为有理函数不恒为零。取坐标 \ensuremath{x_1},\ensuremath{\ldots},\ensuremath{x_{k−1}},T=\ensuremath{x_{k}}，在 K=\ensuremath{F_2}(\ensuremath{x_1},\ensuremath{\ldots},\ensuremath{x_{k−1}}) 中令 L(T)=\ensuremath{\prod}\_\{u\ensuremath{\in}\ensuremath{F_2}\textasciicircum\{\}\{k−1\}\}(T+u·x\ensuremath{\prime})。L 是线性化多项式，其普通导数 L\ensuremath{\prime}(T)=C 为 T 的一次项系数，且 C=\ensuremath{\prod}\_\{0\ensuremath{\ne}u\ensuremath{\in}\ensuremath{F_2}\textasciicircum\{\}\{k−1\}\}(u·x\ensuremath{\prime})\ensuremath{\ne}0。于是

\ensuremath{\Sigma}\_\{u\ensuremath{\in}\ensuremath{F_2}\textasciicircum\{\}\{k−1\}\} 1/(T+u·x\ensuremath{\prime}) = L\ensuremath{\prime}(T)/L(T) = C/L(T).

把 \ensuremath{\ell}\_k 按最后坐标是否为 0 分组得 \ensuremath{\ell}\_k(x\ensuremath{\prime},T)=\ensuremath{\ell}\_\{k−1\}(x\ensuremath{\prime})+C/L(T)。第二项是 T 的非常数有理函数，而第一项与 T 无关，故二者不能恒等抵消。因此 \ensuremath{\ell}\_k\ensuremath{\ne}0。\ensuremath{R_{V}} 是 domain，所以该非零 homogeneous section 是 nonzerodivisor；反转 grading 后 \ensuremath{R_{V}} 由 degree-one reciprocal generators 生成，O(1) 为 ample invertible sheaf，故 \ensuremath{D_{k}} 为 effective ample Cartier divisor。\ensuremath{Q_{V}} 纯维 k−1，故 dim \ensuremath{D_{k}}=k−2。

\end{proof}

\begin{quote}\textbf{命题 3.2（几何点与零和子空间）. 对 q=\ensuremath{2^{n}}，有 \ensuremath{X_{k}}(\ensuremath{F_{q}})\ensuremath{\ne}\ensuremath{\varnothing} 当且仅当 \ensuremath{F_{q}} 中存在 k 维 \ensuremath{F_2}-子空间 E 使 \ensuremath{\Sigma}\_\{0\ensuremath{\ne}x\ensuremath{\in}E\}\ensuremath{x^{-1}}=0。}\end{quote}

\begin{proof}[证明] 

\ensuremath{\Omega}\_k(\ensuremath{F_{q}}) 的点可写成 [\ensuremath{x_1}:\ensuremath{\cdots}:\ensuremath{x_{k}}]，且它避开所有 \ensuremath{F_2}-rational hyperplanes 当且仅当 \ensuremath{x_1},\ensuremath{\ldots},\ensuremath{x_{k}} 在 \ensuremath{F_{q}} 中 \ensuremath{F_2}-线性无关。令 E=\ensuremath{\langle}\ensuremath{x_1},\ensuremath{\ldots},\ensuremath{x_{k}}\ensuremath{\rangle}\_\{\ensuremath{F_2}\}，则 E\ensuremath{\setminus}\{0\} 恰由所有 0\ensuremath{\ne}v\ensuremath{\in}\ensuremath{F_2}\textasciicircum\{\}k 的 v·x 遍历。因此 \ensuremath{\ell}\_k(x)=\ensuremath{\Sigma}\_\{0\ensuremath{\ne}v\}(v·x)\ensuremath{^{-1}}=\ensuremath{\Sigma}\_\{0\ensuremath{\ne}y\ensuremath{\in}E\}\ensuremath{y^{-1}}。项目化缩放只把 E 换成 \ensuremath{\lambda}E，而 reciprocal sum 乘以 \ensuremath{\ensuremath{\lambda}^{-1}}，不改变“是否为零”。故结论成立。

\end{proof}

\section{满秩开层的显式二次完全交模型}

\subsection{递归坐标}

由于 \ensuremath{\Omega}\_r 删除坐标超平面 \ensuremath{x_1}=0，整个 \ensuremath{\Omega}\_r 都位于 \ensuremath{x_1}\ensuremath{\ne}0 的单一 affine chart 中。故可全局归一化 \ensuremath{x_1}=1。令 \ensuremath{E_{j}}=\ensuremath{\langle}\ensuremath{x_1},\ensuremath{\ldots},\ensuremath{x_{j}}\ensuremath{\rangle}，\ensuremath{L_{j}}=\ensuremath{L_{E_j}}，并置 \ensuremath{L_{0}}(X)=X、\ensuremath{\ensuremath{\delta}_1}=1、\ensuremath{\delta}\_j=\ensuremath{L_{j−1}}(\ensuremath{x_{j}})（j\ensuremath{\ge}2）。若 \ensuremath{C_{j}}=[X]\ensuremath{L_{j}}，则引理 2.3 给出

\ensuremath{L_{j}}(X)=\ensuremath{L_{j−1}}(X)²+\ensuremath{\delta}\_jL\_\{j−1\}(X),    \ensuremath{C_{j}}=\ensuremath{\delta}\_jC\_\{j−1\}.    (4.1)

再定义 Pink 的 unipotent-type reciprocal functions

\ensuremath{f_{j}} = \ensuremath{\Sigma}\_\{u\ensuremath{\in}\ensuremath{E_{j−1}}\} 1/(\ensuremath{x_{j}}+u) = \ensuremath{C_{j−1}}/\ensuremath{\delta}\_j.    (4.2)

第二个等号由 \ensuremath{L_{j−1}} 的对数导数得到。由 (4.1) 立刻有

\ensuremath{f_{j}}/\ensuremath{f_{j+1}}=\ensuremath{\delta}\_\{j+1\}/\ensuremath{\delta}\_\ensuremath{j^2}.    (4.3)

此外，\ensuremath{F_2}\textasciicircum\{\}r\ensuremath{\setminus}\{0\} 可按“最后一个非零坐标”分块；第 j 块正对应 \ensuremath{x_{j}}+\ensuremath{E_{j−1}}，故

\ensuremath{\ell}\_r=\ensuremath{f_1}+\ensuremath{f_2}+\ensuremath{\cdots}+\ensuremath{f_{r}},    \ensuremath{f_1}=1.    (4.4)

\subsection{二次 scheme \ensuremath{Y_{r}}}

\begin{quote}\textbf{定义 4.1. 设 \ensuremath{Y_{r}}\ensuremath{\subset}\ensuremath{A^{N_{r}}}\_\{\ensuremath{F_2}\} 由下列变量和方程定义。变量：\ensuremath{u_{a,j}}（2\ensuremath{\le}a\ensuremath{\le}r, 1\ensuremath{\le}j<a），\ensuremath{t_{i}},\ensuremath{d_{i}}（2\ensuremath{\le}i\ensuremath{\le}r），\ensuremath{c_{j}}（1\ensuremath{\le}j<r）；约定 \ensuremath{u_{a,a}}=1、\ensuremath{t_1}=1。方程：(i) \ensuremath{u_{a,j}}²+\ensuremath{u_{a,j}}=\ensuremath{c_{j}} \ensuremath{u_{a,j+1}}；(ii) \ensuremath{t_{i}} \ensuremath{d_{i}}=1；(iii) \ensuremath{c_1}\ensuremath{t_2}=1，c\_jt\_\{j+1\}=\ensuremath{t_{j}}（2\ensuremath{\le}j<r）；(iv) 1+\ensuremath{t_2}+\ensuremath{\cdots}+\ensuremath{t_{r}}=0。}\end{quote}

变量数和方程数分别为

\ensuremath{N_{r}} = C(r,2)+3(r−1),    \ensuremath{R_{r}}=C(r,2)+2(r−1)+1,    (4.5)

\ensuremath{N_{r}}−\ensuremath{R_{r}}=r−2.    (4.6)

\subsection{Scheme-level 同构}

\begin{quote}\textbf{定理 4.2（显式二次完全交模型）. 对每个 r\ensuremath{\ge}2，Carlet 满秩开层 \ensuremath{X_{r}} 与 \ensuremath{Y_{r}} 在 \ensuremath{F_2} 上 scheme-theoretically 同构。  [本文自建关键引理]}\end{quote}

\begin{proof}[证明] 

第一步：构造 \ensuremath{X_{r}}\ensuremath{\to}\ensuremath{Y_{r}}。对 \ensuremath{X_{r}} 上的规范化基 \ensuremath{x_1}=1,\ensuremath{x_2},\ensuremath{\ldots},\ensuremath{x_{r}}，定义

\ensuremath{t_{j}}=\ensuremath{f_{j}},    \ensuremath{d_{j}}=\ensuremath{t_{j}}\ensuremath{^{-1}},    \ensuremath{c_{j}}=\ensuremath{t_{j}}/\ensuremath{t_{j+1}},    \ensuremath{u_{a,j}}=\ensuremath{L_{j−1}}(\ensuremath{x_{a}})/\ensuremath{\delta}\_j.    (4.7)

所有分母均为单位：\ensuremath{\Omega}\_r 的定义保证每个非零 \ensuremath{F_2}-线性组合可逆，因此 \ensuremath{\delta}\_j=\ensuremath{\prod}\_\{u\ensuremath{\in}\ensuremath{E_{j−1}}\}(\ensuremath{x_{j}}+u) 为单位；\ensuremath{C_{j−1}}=\ensuremath{\prod}\_\{0\ensuremath{\ne}u\ensuremath{\in}\ensuremath{E_{j−1}}\}u 亦为单位，故 \ensuremath{f_{j}}=\ensuremath{C_{j−1}}/\ensuremath{\delta}\_j 为单位。于是 (4.7) 给出正则函数。

验证 Artin–Schreier 方程。由 (4.1)，对 a>j 有

\ensuremath{u_{a,j}}²+\ensuremath{u_{a,j}} = \ensuremath{L_{j}}(\ensuremath{x_{a}})/\ensuremath{\delta}\_\ensuremath{j^2} = (\ensuremath{\delta}\_\{j+1\}/\ensuremath{\delta}\_\ensuremath{j^2})\ensuremath{u_{a,j+1}} = \ensuremath{c_{j}} \ensuremath{u_{a,j+1}}.

其中最后一步用 (4.3)。逆元方程和比例方程由定义立即成立；零和方程则由 (4.4) 和 \ensuremath{\ell}\_r=0 得到。因此存在 morphism Φ:\ensuremath{X_{r}}\ensuremath{\to}\ensuremath{Y_{r}}。

第二步：构造 \ensuremath{Y_{r}}\ensuremath{\to}\ensuremath{X_{r}}。令 A 为任意 \ensuremath{F_2}-algebra，并取一个 A-valued point of \ensuremath{Y_{r}}。置 \ensuremath{x_1}=1、\ensuremath{x_{a}}=\ensuremath{u_{a,1}}。递归定义 \ensuremath{\ensuremath{\delta}_1}=1、\ensuremath{\delta}\_\{j+1\}=c\_j\ensuremath{\delta}\_\ensuremath{j^2}。我们声称对所有 a\ensuremath{\ge}j 有

\ensuremath{L_{j−1}}(\ensuremath{x_{a}})=\ensuremath{\delta}\_j \ensuremath{u_{a,j}}.    (4.8)

对 j=1，(4.8) 就是 \ensuremath{x_{a}}=\ensuremath{u_{a,1}}。假设对 j 成立，则由子空间多项式递推和 \ensuremath{Y_{r}} 的二次方程，

\ensuremath{L_{j}}(\ensuremath{x_{a}})=\ensuremath{\delta}\_\ensuremath{j^2}(\ensuremath{u_{a,j}}²+\ensuremath{u_{a,j}})=\ensuremath{\delta}\_\ensuremath{j^2}c\_j \ensuremath{u_{a,j+1}}=\ensuremath{\delta}\_\{j+1\}\ensuremath{u_{a,j+1}},

故归纳成立。取 a=j 得 \ensuremath{L_{j−1}}(\ensuremath{x_{j}})=\ensuremath{\delta}\_j，因为约定 \ensuremath{u_{j,j}}=1。

第三步：证明反向构造真正落在 Drinfeld open。由 \ensuremath{t_{i}} \ensuremath{d_{i}}=1，所有 \ensuremath{t_{i}} 是单位。比例方程给 \ensuremath{c_1}=\ensuremath{t_2}\ensuremath{^{-1}} 以及 \ensuremath{c_{j}}=t\_jt\_\{j+1\}\ensuremath{^{-1}}，故所有 \ensuremath{c_{j}} 是单位，继而所有 \ensuremath{\delta}\_j 是单位。由于

\ensuremath{\delta}\_j=\ensuremath{L_{j−1}}(\ensuremath{x_{j}})=\ensuremath{\prod}\_\{u\ensuremath{\in}\ensuremath{E_{j−1}}\}(\ensuremath{x_{j}}+u),

在任意交换环中乘积为单位会强迫每个因子为单位，因此所有 \ensuremath{x_{j}}+u（u\ensuremath{\in}\ensuremath{E_{j−1}}）都是单位。任一非零 \ensuremath{F_2}-线性组合取其最大非零指标 j，都可写成 \ensuremath{x_{j}}+u；故所有非零线性组合为单位。这正是 \ensuremath{\Omega}\_r 的 functor-of-points 条件，故反向 morphism Ψ:\ensuremath{Y_{r}}\ensuremath{\to}\ensuremath{\Omega}\_r 合法。

第四步：证明反向点满足 Carlet 方程。设 \ensuremath{C_{j}}=[X]\ensuremath{L_{j}}。由 \ensuremath{C_{j}}=\ensuremath{\delta}\_jC\_\{j−1\} 和 (4.2)，反向构造的实际 reciprocal functions 满足

\ensuremath{f_{j}}/\ensuremath{f_{j+1}}=\ensuremath{\delta}\_\{j+1\}/\ensuremath{\delta}\_\ensuremath{j^2}=\ensuremath{c_{j}}=\ensuremath{t_{j}}/\ensuremath{t_{j+1}}.

又 \ensuremath{f_1}=\ensuremath{t_1}=1，故归纳得 \ensuremath{f_{j}}=\ensuremath{t_{j}} 对所有 j 成立。\ensuremath{Y_{r}} 的最后方程 1+\ensuremath{t_2}+\ensuremath{\cdots}+\ensuremath{t_{r}}=0 因而等价于 \ensuremath{\ell}\_r=\ensuremath{f_1}+\ensuremath{\cdots}+\ensuremath{f_{r}}=0，所以 Ψ 的像落在 \ensuremath{X_{r}}。

第五步：验证互逆。从 \ensuremath{X_{r}} 出发，反向映射恢复 \ensuremath{x_{a}}=\ensuremath{u_{a,1}}=\ensuremath{L_{0}}(\ensuremath{x_{a}})/\ensuremath{\ensuremath{\delta}_1}=\ensuremath{x_{a}}；从 \ensuremath{Y_{r}} 出发，(4.8) 保证正向重新计算的 \ensuremath{u_{a,j}} 等于原 \ensuremath{u_{a,j}}，而上一段保证 \ensuremath{t_{j}}=\ensuremath{f_{j}}，\ensuremath{d_{j}}=\ensuremath{t_{j}}\ensuremath{^{-1}}、\ensuremath{c_{j}}=\ensuremath{t_{j}}/\ensuremath{t_{j+1}} 也全部恢复。因此 Φ\ensuremath{\circ}Ψ 与 Ψ\ensuremath{\circ}Φ 在所有坐标环生成元上均为恒等，故 Φ、Ψ 是 scheme isomorphisms。

\end{proof}

\subsection{光滑性与有限 étale 结构}

\begin{quote}\textbf{推论 4.3. \ensuremath{X_{r}} 是光滑纯维 r−2 的 affine \ensuremath{F_2}-scheme。更精确地，令 \ensuremath{T_{r}}=\{(\ensuremath{t_2},\ensuremath{\ldots},\ensuremath{t_{r}})\ensuremath{\in}\ensuremath{G_{m}}\textasciicircum\{\}\{r−1\}:1+\ensuremath{t_2}+\ensuremath{\cdots}+\ensuremath{t_{r}}=0\}，则 \ensuremath{X_{r}}\ensuremath{\cong}\ensuremath{Y_{r}}\ensuremath{\to}\ensuremath{T_{r}} 为次数 \ensuremath{2^{C(r,2)}} 的 finite étale morphism。}\end{quote}

\begin{proof}[证明] 

由 \ensuremath{t_{i}} \ensuremath{d_{i}}=1 与比例方程可消去全部 \ensuremath{d_{i}}、\ensuremath{c_{j}}，因此底空间就是 \ensuremath{T_{r}}。它是 torus 中一个线性超平面，偏导数中任一个均为 1，故光滑维 r−2。固定 t 后，对每个 \ensuremath{u_{a,j}}，方程 \ensuremath{u_{a,j}}²+\ensuremath{u_{a,j}}=\ensuremath{c_{j}} \ensuremath{u_{a,j+1}} 对当前变量是 monic 二次多项式，且导数恒为 1。按 j 从 a−1 向 1 递归，它给出迭代 finite étale 二次扩张。共有 \ensuremath{\Sigma}\_\{a=2\}\textasciicircum\{\}r(a−1)=C(r,2) 步，故次数为 \ensuremath{2^{C(r,2)}}。

\end{proof}

\section{秩分层与 rational smoothness}

\paragraph{局部交叉上同调的常值层识别.} 对 reciprocal cone 顶点的局部交叉上同调计算，应把模格的模性结论解释为局部 intersection complex 与常值层的自然识别；由于构造定义在 $\mathbb F_2$ 上，该识别与 Frobenius 作用相容。与光滑因子作乘积时同时加入相应的上同调位移。

\subsection{Pink–Schieder rank stratification}

Pink–Schieder [PS, Thm. 8.2] 证明 \ensuremath{Q_{V}} 是所有非零子空间 W\ensuremath{\le}V 所对应 \ensuremath{\Omega}\_W 的 set-theoretic disjoint union，且 closure relation 为 \ensuremath{\Omega}̄\_W=\ensuremath{\bigcup}\_\{0\ensuremath{\ne}W\ensuremath{\prime}\ensuremath{\le}W\}\ensuremath{\Omega}\_\{W\ensuremath{\prime}\}。其 extension-by-zero 的模解释还表明 reciprocal coordinate \ensuremath{z_{v}} 在 \ensuremath{\Omega}\_W 上若 v\ensuremath{\notin}W 则结构层意义下拉回为 0。

\begin{quote}\textbf{命题 5.1. Carlet divisor 具有分层}\end{quote}

\ensuremath{D_{k}} = \ensuremath{\bigsqcup}\_\{r=2\}\textasciicircum\{\}k  \ensuremath{\bigsqcup}\_\{W\ensuremath{\le}\ensuremath{F_2}\textasciicircum\{\}k, dim W=r\} \ensuremath{X_{r}}.    (5.1)

\begin{proof}[证明] 

在 \ensuremath{\Omega}\_W 上，extension-by-zero 使 \ensuremath{\ell}\_k 的限制精确为 \ensuremath{\ell}\_W，因此 \ensuremath{D_{k}}\ensuremath{\cap}\ensuremath{\Omega}\_W=\ensuremath{X_{dim W}}。rank 1 时 \ensuremath{\ell}\_1=1/x\ensuremath{\ne}0，故无 rank-1 点。不同 W 的 strata 由 [PS, Thm. 8.2] 不交，得 (5.1)。闭包关系同时给出按 rank 的闭过滤。

\end{proof}

\subsection{Reciprocal normal cone 的 local IH}

Pink–Schieder [PS, Prop. 8.3] 给出更强的局部乘积。若 V=W\ensuremath{\oplus}V\ensuremath{\prime\prime}，则在 \ensuremath{\Omega}\_W 的某个 Zariski 邻域中，\ensuremath{Q_{V}} 与 \ensuremath{\Omega}\_W\ensuremath{\times}Q̃\_\{V\ensuremath{\prime\prime}\} 在 \ensuremath{\Omega}\_W\ensuremath{\times}\{0\} 的邻域同构，其中 Q̃\_\{V\ensuremath{\prime\prime}\}=Spec \ensuremath{R_{V\prime\prime}} 是 affine reciprocal cone。

\begin{quote}\textbf{引理 5.2. affine reciprocal cone Q̃\_U 在顶点 0 处具有平凡的 \ensuremath{\ell}-adic local intersection cohomology（\ensuremath{\ell}\ensuremath{\ne}2）。}\end{quote}

\begin{proof}[证明] 

把 U\ensuremath{\setminus}\{0\} 看作 U* 上的一组线性形式，\ensuremath{R_{U}} 正是相应 Orlik–Terao/reciprocal-plane coordinate ring。其 matroid 是二元 projective geometry PG(dim U−1,2)，flat lattice 等同于 U 的子空间格。子空间格是 modular lattice，且其任意 interval 仍 modular。

Elias–Proudfoot–Wakefield [EPW, Prop. 2.14] 证明对 matroid M，\ensuremath{P_{M}}(t)=1 当且仅当 L(M) modular；[EPW, Thm. 3.10] 将 representable matroid 的 KL polynomial 解释为 reciprocal plane 的 \ensuremath{\ell}-adic intersection-cohomology Poincaré polynomial，而 [EPW, Lem. 3.8] 及其 cone-point contraction 论证把顶点 local IH 化为该 reciprocal slice 的 IH。因此这里的 local IH 只有光滑点应有的一维项。

\end{proof}

\subsection{\ensuremath{Q_{k}} 的 rational smoothness}

\begin{quote}\textbf{命题 5.3. \ensuremath{Q_{k}} 是 \ensuremath{\ell}-adically rationally smooth。  [本文组合论证]}\end{quote}

\begin{proof}[证明] 

取任意 x\ensuremath{\in}\ensuremath{\Omega}\_W。由 [PS, Prop. 8.3]，x 的 Zariski 邻域在 x 附近同构于 \ensuremath{\Omega}\_W\ensuremath{\times}Q̃\_\{V\ensuremath{\prime\prime}\} 在 (x,0) 附近。\ensuremath{\Omega}\_W 光滑，而引理 5.2 给 normal factor 顶点平凡 local IH。intersection complex 对光滑乘积的局部行为给出 x 处与同维 affine space 相同的 local IH。因此 \ensuremath{Q_{k}} 每一点 rationally smooth。

\end{proof}

\subsection{Carlet divisor 的 rational smoothness}

\begin{quote}\textbf{引理 5.4（smooth slicing）. 设 p:U\ensuremath{\to}N 在 x 处光滑，相对维 d，h\ensuremath{\in}Γ(U,\ensuremath{O_{U}})，且 dh 在 p 的 fiber tangent direction 上于 x 非零。则 V(h)\ensuremath{\to}N 在 x 处光滑相对维 d−1；因此 étale-locally，(V(h),x)\ensuremath{\cong}(N,p(x))\ensuremath{\times}\ensuremath{A^{d−1}}。}\end{quote}

\begin{proof}[证明] 

在 x 处，非零的相对微分意味着 (p,h):U\ensuremath{\to}N\ensuremath{\times}\ensuremath{A^1} 在 x 处光滑，相对维 d−1（Jacobian criterion / smoothness criterion）。取 \ensuremath{A^1} 中 0 的 fiber 即得 V(h)\ensuremath{\to}N 光滑。任意 smooth morphism étale-locally 是 affine-space projection，故第二个断言成立。

\end{proof}

\begin{quote}\textbf{定理 5.5. \ensuremath{D_{k}} 是 \ensuremath{\ell}-adically rationally smooth。  [本文自建关键定理]}\end{quote}

\begin{proof}[证明] 

取 x\ensuremath{\in}\ensuremath{D_{k}}，位于某个 rank-r stratum \ensuremath{X_{r}}\ensuremath{\subset}\ensuremath{\Omega}\_W。选补空间 V\ensuremath{\prime\prime}，由 [PS, Prop. 8.3] 在 x 附近把 \ensuremath{Q_{k}} 写成 \ensuremath{\Omega}\_W\ensuremath{\times}N，其中 N=Q̃\_\{V\ensuremath{\prime\prime}\}，x 对应 (x,0)。由 extension-by-zero，Carlet section 在 stratum 上的限制是 \ensuremath{\ell}\_W。

由推论 4.3，\ensuremath{X_{r}}=\{\ensuremath{\ell}\_W=0\}\ensuremath{\subset}\ensuremath{\Omega}\_W 是光滑 codimension-one 子空间，因此 d\ensuremath{\ell}\_W 在 x 的 \ensuremath{\Omega}\_W-tangent direction 上非零。把局部 Carlet section 记为 h，并对投影 \ensuremath{\Omega}\_W\ensuremath{\times}N\ensuremath{\to}N 应用引理 5.4，得到 étale-locally

(\ensuremath{D_{k}},x) \ensuremath{\cong}\_ét (N,0) \ensuremath{\times} \ensuremath{A^{r−2}}.    (5.2)

\begin{quote}\textbf{引理 5.2 说明 N 在 0 处 local IH 平凡；乘以光滑 affine factor 只产生标准 shift。故 \ensuremath{D_{k}} 在 x 处 local IH 与光滑 m=k−2 维空间相同。x 任意，故 \ensuremath{D_{k}} rationally smooth。}\end{quote}

\end{proof}

\section{上同调：universal homeomorphism、weak Lefschetz 与对偶}

\paragraph{对偶中的 Tate twist.} 对光滑且有理光滑的 proper $m$ 维对象，Verdier 对偶给出 Frobenius-equivariant 的 Poincar\'e 型对偶 $H^i\cong H^{2m-i,\vee}(-m)$；在迹公式中保留该 Tate twist，才能逐项确定 Frobenius 特征值。

\subsection{Ambient \ensuremath{Q_{k}} 的 Frobenius-equivariant cohomology}

\begin{quote}\textbf{命题 6.1. 对任意 \ensuremath{\ell}\ensuremath{\ne}2，\ensuremath{Q_{k}} 与 \ensuremath{P^{k−1}} 的 \ensuremath{\ell}-adic étale cohomology 作为 Frobenius representations 同构。}\end{quote}

\begin{proof}[证明] 

Pink–Schieder [PS, Thm. 9.3] 构造 \ensuremath{Q_{V}}\ensuremath{\to}P(V*) 与 P(V*)\ensuremath{\to}\ensuremath{Q_{V}}，两方向复合均为 \ensuremath{2^{k−1}}-power Frobenius，并明确指出二者 finite、underlying-set bijective、radicial。因此它们是 finite universal homeomorphisms。

由 étale site 的 topological invariance（例如 Stacks Project, Prop. 59.45.4, Tag 03SI），universal homeomorphism 诱导 étale topoi 与 étale cohomology 的等价。上述 morphisms 定义在 \ensuremath{F_2} 上，pullback 与几何 Frobenius 自然交换，故同构为 Frobenius-equivariant。P(V*)\ensuremath{\cong}\ensuremath{P^{k−1}}，结论成立。

\end{proof}

\subsection{Weak Lefschetz}

置 d=dim \ensuremath{Q_{k}}=k−1，m=dim \ensuremath{D_{k}}=k−2=d−1。由命题 5.3、定理 5.5，

IC\_\{\ensuremath{Q_{k}}\}\ensuremath{\cong}Q\_\ensuremath{\ell}[d],    IC\_\{\ensuremath{D_{k}}\}\ensuremath{\cong}Q\_\ensuremath{\ell}[m].    (6.1)

Wan–Zhang [WZ, arXiv:2501.12623, Lem. 2.3.1] 的 weak Lefschetz 形式如下：若 Y projective、F\ensuremath{\in}\textasciicircum\{\}\ensuremath{pD^{\ge0}}(Y)，X 是 r 个 ample Cartier divisors 的交，则 \ensuremath{H^{j−r}}(Y;F)\ensuremath{\to}\ensuremath{H^{j−r}}(X;F) 在 j<0 时同构、j=0 时单射。我们只用 r=1。

\begin{quote}\textbf{定理 6.2. 自然限制映射 \ensuremath{H^{i}}(\ensuremath{Q_{k}},Q\_\ensuremath{\ell})\ensuremath{\to}\ensuremath{H^{i}}(\ensuremath{D_{k}},Q\_\ensuremath{\ell}) 在 i<m 时为同构，在 i=m 时为单射。}\end{quote}

\begin{proof}[证明] 

\ensuremath{D_{k}} 是 ample Cartier divisor（引理 3.1）。取 F=IC\_\{\ensuremath{Q_{k}}\}=Q\_\ensuremath{\ell}[d]，这是 perverse sheaf，满足 Wan–Zhang 的 cosupport 条件。其 weak Lefschetz 映射写成 \ensuremath{H^{j−1}}(\ensuremath{Q_{k}};Q\_\ensuremath{\ell}[d])\ensuremath{\to}\ensuremath{H^{j−1}}(\ensuremath{D_{k}};Q\_\ensuremath{\ell}[d])。两边分别等于普通上同调 degree i=j−1+d=j+m。故 j<0 正好对应 i<m，j=0 对应 i=m。

\end{proof}

\subsection{高次数由 \ensuremath{D_{k}} 自身对偶恢复}

\begin{quote}\textbf{定理 6.3. 对所有 i\ensuremath{\ne}m，有 Frobenius-equivariant 同构 \ensuremath{H^{i}}(\ensuremath{D_{k}},Q\_\ensuremath{\ell})\ensuremath{\cong}\ensuremath{H^{i}}(\ensuremath{P^{m}},Q\_\ensuremath{\ell})。}\end{quote}

\begin{proof}[证明] 

i<m 时，由定理 6.2 与命题 6.1，\ensuremath{H^{i}}(\ensuremath{D_{k}})\ensuremath{\cong}\ensuremath{H^{i}}(\ensuremath{Q_{k}})\ensuremath{\cong}\ensuremath{H^{i}}(\ensuremath{P^{m+1}})；在 i<m 的范围内，这与 \ensuremath{P^{m}} 的同次数上同调一致。

\ensuremath{D_{k}} proper 且 rationally smooth，故其 intersection complex 是 shifted constant sheaf，并由 Verdier self-duality 得到 ordinary cohomology 的 Poincaré-type duality

\ensuremath{H^{i}}(\ensuremath{D_{k}}) \ensuremath{\cong} \ensuremath{H^{2m−i}}(\ensuremath{D_{k}})\textasciicircum\{\}\ensuremath{\lor}(−m).    (6.2)

若 i>m，则 2m−i<m，右侧已经由前一段算清，并给出相应 Tate twist。对 \ensuremath{P^{m}} 也有相同对偶。因此 \ensuremath{H^{i}}(\ensuremath{D_{k}})\ensuremath{\cong}\ensuremath{H^{i}}(\ensuremath{P^{m}})。Frobenius-equivariance 由各自然构造得到。

\end{proof}

\subsection{点数公式}

在 middle degree，定理 6.2 仅给单射。定义 Frobenius-stable quotient

\ensuremath{V_{k}} := \ensuremath{H^{m}}(\ensuremath{D_{k}},Q\_\ensuremath{\ell}) / im[\ensuremath{H^{m}}(\ensuremath{Q_{k}},Q\_\ensuremath{\ell})\ensuremath{\to}\ensuremath{H^{m}}(\ensuremath{D_{k}},Q\_\ensuremath{\ell})],    \ensuremath{b_{k}}=dim \ensuremath{V_{k}}.    (6.3)

\begin{quote}\textbf{推论 6.4. 对 q=\ensuremath{2^{n}}，}\end{quote}

\#\ensuremath{D_{k}}(\ensuremath{F_{q}})=1+q+\ensuremath{\cdots}+\ensuremath{q^{m}}+(−1)\textasciicircum\{\}m Tr(\ensuremath{F_{q}}|\ensuremath{V_{k}}).    (6.4)

\begin{proof}[证明] 

除 middle degree 外，定理 6.3 给出与 \ensuremath{P^{m}} 完全相同的 Tate cohomology。若 m 偶，则 \ensuremath{H^{m}}(\ensuremath{Q_{k}})=Q\_\ensuremath{\ell}(−m/2) 通过定理 6.2 注入 middle degree，正好提供 \ensuremath{P^{m}} 的 middle Tate 类 \ensuremath{q^{m/2}}；若 m 奇则该项为 0。把剩余 quotient 记为 \ensuremath{V_{k}}，对 proper \ensuremath{D_{k}} 应用 Grothendieck–Lefschetz trace formula 即得。

\end{proof}

\begin{quote}\textbf{引理 6.5. |Tr(\ensuremath{F_{q}}|\ensuremath{V_{k}})|\ensuremath{\le}\ensuremath{b_{k}} \ensuremath{q^{m/2}}。}\end{quote}

\begin{proof}[证明] 

不需要假设 \ensuremath{D_{k}} geometrically irreducible 或 middle cohomology pure。Deligne 的一般权重定理给 \ensuremath{H_{c}}\textasciicircum\{\}m 的 Frobenius weights \ensuremath{\le}m；\ensuremath{D_{k}} proper，故 \ensuremath{H_{c}}\textasciicircum\{\}m=\ensuremath{H^{m}}。\ensuremath{V_{k}} 是 Frobenius-stable quotient，所有 eigenvalues 仍满足绝对值 \ensuremath{\le}\ensuremath{q^{m/2}}，于是 trace 的绝对值至多 b\_kq\textasciicircum\{\}\{m/2\}。

\end{proof}

\section{显式 middle Betti 上界}

\paragraph{紧支撑上同调的次数范围.} 这里 $X_r$ 光滑、仿射、纯维 $d=r-2$。Artin 消失给出 $H^j(X_r)=0$（$j>d$），Poincar\'e 对偶遂给出 $H_c^i(X_r)=0$（$i<d$）；维数界给出 $H_c^i(X_r)=0$（$i>2d$）。因此下式确实覆盖总紧支撑 Betti 数的全部可能次数；$r=2$ 时 $d=0$ 亦成立。

\subsection{Wan–Zhang affine complete-intersection bound}

我们使用 Wan–Zhang [WZ, Cor. 4.2.4(ii)]：若 V\ensuremath{\subset}\ensuremath{A^{N}} 由 R 个次数 \ensuremath{\le}d 的方程定义且 dim V=N−R，则对 0\ensuremath{\le}j\ensuremath{\le}N−R，

dim \ensuremath{H_{c}}\textasciicircum\{\}\{N−R+j\}(V,Q\_\ensuremath{\ell}) \ensuremath{\le} C(N−j−1,R−1)\ensuremath{d^{N−j}}.    (7.1)

由定理 4.2，\ensuremath{X_{r}}\ensuremath{\cong}\ensuremath{Y_{r}}；由推论 4.3，dim \ensuremath{X_{r}}=r−2=\ensuremath{N_{r}}−\ensuremath{R_{r}}；所有 defining equations 次数 \ensuremath{\le}2。因此

dim \ensuremath{H_{c}}\textasciicircum\{\}\{r−2+j\}(\ensuremath{X_{r}}) \ensuremath{\le} C(\ensuremath{N_{r}}−j−1,\ensuremath{R_{r}}−1)\ensuremath{2^{N_{r−j}}},  0\ensuremath{\le}j\ensuremath{\le}r−2.    (7.2)

\begin{quote}\textbf{引理 7.1. 总 compact-support Betti 数满足 \ensuremath{B_{c}}(\ensuremath{X_{r}})<\ensuremath{A_{r}}，其中}\end{quote}

\ensuremath{A_{r}} := \ensuremath{2^{N_{r+1}}} C(\ensuremath{N_{r}}−1,r−2).    (7.3)

\begin{proof}[证明] 

因 \ensuremath{R_{r}}=\ensuremath{N_{r}}−(r−2)，有 C(\ensuremath{N_{r}}−j−1,\ensuremath{R_{r}}−1)=C(\ensuremath{N_{r}}−j−1,r−2−j)\ensuremath{\le}C(\ensuremath{N_{r}}−1,r−2)。式 (7.2) 对 j=0,\ensuremath{\ldots},r−2 求和，并用 \ensuremath{\Sigma}\_\{j\ensuremath{\ge}0\}\ensuremath{2^{−j}}<2，得到 (7.3)。

\end{proof}

\subsection{Rank filtration}

由命题 5.1 及 [PS, Thm. 8.2] 的 closure relation，按 rank\ensuremath{\le}r 的 strata 之并给出 \ensuremath{D_{k}} 的闭过滤。associated compact-support spectral sequence 的 \ensuremath{E_1} 页是各 strata 的 \ensuremath{H_{c}}。记 \ensuremath{G_2}(k,r) 为二元 Gaussian binomial coefficient，即 V 中 r 维子空间的个数。因此

B(\ensuremath{D_{k}})=\ensuremath{B_{c}}(\ensuremath{D_{k}}) \ensuremath{\le} \ensuremath{\Sigma}\_\{r=2\}\textasciicircum\{\}k \ensuremath{G_2}(k,r) \ensuremath{B_{c}}(\ensuremath{X_{r}}).    (7.4)

\begin{quote}\textbf{引理 7.2（Gaussian bound）. 对 0\ensuremath{\le}r\ensuremath{\le}k，有 \ensuremath{G_2}(k,r) < 4·\ensuremath{2^{r(k−r)}}。}\end{quote}

\begin{proof}[证明] 

写 \ensuremath{G_2}(k,r)=\ensuremath{2^{r(k−r)}}\ensuremath{\prod}\_\{i=0\}\textasciicircum\{\}\{r−1\}(1−\ensuremath{2^{i−k}})/(1−\ensuremath{2^{i−r}})。分子 <1，分母倒数的无限乘积 \ensuremath{\prod}\_\{j\ensuremath{\ge}1\}(1−\ensuremath{2^{-j}})\textasciicircum\{\}\{-1\}<4，故结论。

\end{proof}

\begin{quote}\textbf{引理 7.3. 若 r=k−s，则}\end{quote}

\ensuremath{N_{k}}−\ensuremath{N_{k−s}}−(k−s)s = s(s+5)/2.    (7.5)

\begin{proof}[证明] 

由 \ensuremath{N_{t}}=t(t−1)/2+3(t−1)=(\ensuremath{t^2}+5t−6)/2 直接展开即可。

\end{proof}

\begin{quote}\textbf{定理 7.4（closed Betti bound）. 对 k\ensuremath{\ge}3，}\end{quote}

\ensuremath{b_{k}} < \ensuremath{\beta}\_k := \ensuremath{2^{N_{k+2}}} C(\ensuremath{N_{k}}−1,k−2).    (7.6)

\begin{proof}[证明] 

由 \ensuremath{b_{k}}\ensuremath{\le}B(\ensuremath{D_{k}})、(7.4) 和引理 7.1，最高 rank r=k 的贡献小于 \ensuremath{A_{k}}。对 r=k−s 的低秩项，先注意 C(\ensuremath{N_{k−s}}−1,k−s−2)\ensuremath{\le}C(\ensuremath{N_{k}}−1,k−2)：先增大上指标，再利用 b\ensuremath{\mapsto}C(\ensuremath{N_{k}}−1,b) 在 0\ensuremath{\le}b\ensuremath{\le}k−2 范围递增；这里 k−2<(\ensuremath{N_{k}}−1)/2。

结合 Gaussian bound、(7.5) 与该二项式比较，低秩项相对于 \ensuremath{A_{k}} 的比例小于

4·2\textasciicircum\{\}\{(k−s)s+\ensuremath{N_{k−s}}−\ensuremath{N_{k}}\} = 4·\ensuremath{2^{−s(s+5)/2}}.

故所有 s\ensuremath{\ge}1 的低秩贡献总比例小于 4\ensuremath{\Sigma}\_\{s\ensuremath{\ge}1\}\ensuremath{2^{−s(s+5)/2}}=0.53224194\ensuremath{\ldots}<1。因此 B(\ensuremath{D_{k}})<2A\_k=\ensuremath{2^{N_{k+2}}}C(\ensuremath{N_{k}}−1,k−2)，从而 \ensuremath{b_{k}} 满足 (7.6)。

\end{proof}

\section{低秩边界的点数上界}

\subsection{清分母多项式}

在普通 projective coordinates x=(\ensuremath{x_1},\ensuremath{\ldots},\ensuremath{x_{r}}) 上定义 Moore-type hyperplane product

Δ\_r(x)=\ensuremath{\prod}\_\{0\ensuremath{\ne}v\ensuremath{\in}\ensuremath{F_2}\textasciicircum\{\}r\}(v·x),    (8.1)

以及

\ensuremath{P_{r}}(x)=Δ\_r(x)\ensuremath{\ell}\_r(x)=\ensuremath{\Sigma}\_\{0\ensuremath{\ne}v\ensuremath{\in}\ensuremath{F_2}\textasciicircum\{\}r\} \ensuremath{\prod}\_\{0\ensuremath{\ne}w\ensuremath{\in}\ensuremath{F_2}\textasciicircum\{\}r, w\ensuremath{\ne}v\}(w·x).    (8.2)

\begin{quote}\textbf{引理 8.1. \ensuremath{P_{r}} 是非零 homogeneous polynomial，deg \ensuremath{P_{r}}=\ensuremath{2^{r}}−2，且 \ensuremath{X_{r}}\ensuremath{\subset}V(\ensuremath{P_{r}})\ensuremath{\subset}\ensuremath{P^{r−1}}。}\end{quote}

\begin{proof}[证明] 

齐次次数由 (8.2) 直接得到。\ensuremath{X_{r}} 位于 Δ\_r\ensuremath{\ne}0 的 Drinfeld open，在该开集 \ensuremath{P_{r}}=Δ\_r\ensuremath{\ell}\_r，故 \ensuremath{\ell}\_r=0 蕴含 \ensuremath{P_{r}}=0。只需证明 \ensuremath{P_{r}}\ensuremath{\ne}0。

写 \ensuremath{x_{r}}=T，其余为 x\ensuremath{\prime}，令 L(T)=\ensuremath{\prod}\_\{u\ensuremath{\in}\ensuremath{F_2}\textasciicircum\{\}\{r−1\}\}(T+u·x\ensuremath{\prime})。与引理 3.1 相同，L\ensuremath{\prime}(T)=C\ensuremath{\ne}0，且

\ensuremath{\ell}\_r(x\ensuremath{\prime},T)=\ensuremath{\ell}\_\{r−1\}(x\ensuremath{\prime})+C/L(T).

右侧包含一个真正依赖 T 的有理函数，故 \ensuremath{\ell}\_r 不是零函数；Δ\_r 也非零，因此 \ensuremath{P_{r}}=Δ\_r\ensuremath{\ell}\_r 非零。

\end{proof}

\subsection{Projective point bound}

\begin{quote}\textbf{引理 8.2. 对任意 q=\ensuremath{2^{n}}，}\end{quote}

\#\ensuremath{X_{r}}(\ensuremath{F_{q}}) < \ensuremath{2^{r}} \ensuremath{q^{r−1}}/(q−1).    (8.3)

\begin{proof}[证明] 

Schwartz–Zippel 对非零 degree d=\ensuremath{2^{r}}−2 的 \ensuremath{P_{r}} 在 \ensuremath{A^{r}}(\ensuremath{F_{q}}) 上给至多 d \ensuremath{q^{r−1}} 个零点。每个 projective zero 有 q−1 个非零 cone representatives，因此 \#V(\ensuremath{P_{r}})(\ensuremath{F_{q}})<\ensuremath{2^{r}} \ensuremath{q^{r−1}}/(q−1)。\ensuremath{X_{r}}\ensuremath{\subset}V(\ensuremath{P_{r}}) 给结论。

\end{proof}

\subsection{Boundary ratio}

由命题 5.1，

\ensuremath{D_{k}} \ensuremath{\setminus} \ensuremath{X_{k}} = \ensuremath{\bigsqcup}\_\{r=2\}\textasciicircum\{\}\{k−1\} \ensuremath{\bigsqcup}\_\{dim W=r\} \ensuremath{X_{r}}.    (8.4)

\begin{quote}\textbf{定理 8.3. 若 k\ensuremath{\ge}17 且 n\ensuremath{\ge}2k+1，q=\ensuremath{2^{n}}，则}\end{quote}

\#(\ensuremath{D_{k}} \ensuremath{\setminus} \ensuremath{X_{k}})(\ensuremath{F_{q}}) / \ensuremath{q^{k−2}} < 1/4 + \ensuremath{2^{−k−5}}.    (8.5)

\begin{proof}[证明] 

由 (8.3) 及 Gaussian symmetry \ensuremath{G_2}(k,r)=\ensuremath{G_2}(k,k−r)，并置 s=k−r，得到

\#boundary/\ensuremath{q^{k−2}} < [q/(q−1)] \ensuremath{\Sigma}\_\{s=1\}\textasciicircum\{\}\{k−2\} \ensuremath{G_2}(k,s) \ensuremath{2^{k−s}} \ensuremath{q^{−s}}.    (8.6)

s=1 时 \ensuremath{G_2}(k,1)=\ensuremath{2^{k}}−1，且 q\ensuremath{\ge}\ensuremath{2^{2k+1}}，故该项为 ((\ensuremath{2^{k}}−1)\ensuremath{2^{k−1}})/(q−1)<1/4。

s\ensuremath{\ge}2 时用引理 7.2 和 q\ensuremath{\ge}\ensuremath{2^{2k+1}}，第 s 项（除共同 q/(q−1) 外）小于 \ensuremath{2^{2+k−s(k+s+2)}}。s=2 的首项为 \ensuremath{2^{−k−6}}；连续两项之比至多 \ensuremath{2^{−(k+7)}}。又 q/(q−1)<1+\ensuremath{2^{-35}}，因此整个 s\ensuremath{\ge}2 尾和小于 \ensuremath{2^{−k−5}}。合并即得。

\end{proof}

\section{大参数正性}

置 m=k−2。由推论 6.4、引理 6.5、定理 7.4 与 boundary identity，

\#\ensuremath{X_{k}}(\ensuremath{F_{q}}) \ensuremath{\ge} \ensuremath{q^{m}} − \ensuremath{\beta}\_k \ensuremath{q^{m/2}} − \#boundary.    (9.1)

我们故意丢掉 1+q+\ensuremath{\cdots}+\ensuremath{q^{m−1}} 这些正项，因此只要右端仍正即可。

\subsection{k=17}

\begin{quote}\textbf{命题 9.1. 若 k=17 且 n\ensuremath{\ge}35，则 X\_17(\ensuremath{\mathbb F_{2^n}})\ensuremath{\ne}\ensuremath{\varnothing}。}\end{quote}

\begin{proof}[证明] 

N\_17=184 且 C(183,15)<\ensuremath{2^{72}}，所以 \ensuremath{\beta}\_17<\ensuremath{2^{184+2+72}}=\ensuremath{2^{258}}。n\ensuremath{\ge}35、m=15 时 \ensuremath{q^{m/2}}\ensuremath{\ge}\ensuremath{2^{262.5}}，因此 middle error ratio \ensuremath{\beta}\_17/\ensuremath{q^{m/2}}<\ensuremath{2^{-4.5}}<1/16。由定理 8.3，boundary ratio<1/4+\ensuremath{2^{-22}}。故 (9.1) 的归一化右端大于 1−1/16−1/4−\ensuremath{2^{-22}}>0。

\end{proof}

\subsection{k at least 18}

\begin{quote}\textbf{命题 9.2. 若 k\ensuremath{\ge}18 且 n\ensuremath{\ge}2k+1，则 \ensuremath{X_{k}}(\ensuremath{\mathbb F_{2^n}})\ensuremath{\ne}\ensuremath{\varnothing}。}\end{quote}

\begin{proof}[证明] 

由标准二项式估计 C(M,r)\ensuremath{\le}(eM/r)\textasciicircum\{\}r，且

(\ensuremath{N_{k}}−1)/(k−2) = (k+7)/2 + 3/(k−2),

对 k\ensuremath{\ge}18 有 e(\ensuremath{N_{k}}−1)/(k−2)<2k。因此 C(\ensuremath{N_{k}}−1,k−2)<(2k)\textasciicircum\{\}\{k−2\}。再用 lo\ensuremath{g_2}k\ensuremath{\le}k/4（k\ensuremath{\ge}18），在最坏 n=2k+1 下

lo\ensuremath{g_2}[\ensuremath{\beta}\_k/\ensuremath{q^{(k−2)/2}}] \ensuremath{\le} −\ensuremath{k^2}/4 + (9/2)k − 2.    (9.2)

右端在 k=18 等于 −2，并对 k\ensuremath{\ge}18 单调递减，故 middle ratio\ensuremath{\le}1/4。boundary ratio 由定理 8.3 小于 1/4+\ensuremath{2^{-k−5}}。于是总误差比例 <1/2+\ensuremath{2^{-23}}<1，(9.1) 给 \#\ensuremath{X_{k}}>0。

\end{proof}

\subsection{临界 k=16}

\begin{quote}\textbf{命题 9.3. 若 k=16 且 n\ensuremath{\ge}37，则 X\_16(\ensuremath{\mathbb F_{2^n}})\ensuremath{\ne}\ensuremath{\varnothing}。}\end{quote}

\begin{proof}[证明] 

N\_16=165 且 C(164,14)<\ensuremath{2^{66}}，所以 \ensuremath{\beta}\_16<\ensuremath{2^{233}}。m=14、n\ensuremath{\ge}37 时 \ensuremath{q^{m/2}}=\ensuremath{q^{7}}\ensuremath{\ge}\ensuremath{2^{259}}，middle ratio<\ensuremath{2^{-26}}。对 boundary 使用 (8.6) 的同样估计（事实上 n=37 比形式阈值 33 更有余量），粗略界 <1/4+\ensuremath{2^{-21}} 已足够。因此总误差远小于 1，故 X\_16 非空。

\end{proof}

\section{有限尾部与候选主定理}

\subsection{外部已知范围}

本文使用以下公开结果作为有限尾部输入：

\begin{itemize}

\item Hou–Zhao [HZ1, Thm. 5.4]：对 n\ensuremath{\ge}8，低半区至少覆盖 3\ensuremath{\le}k\ensuremath{\le}ceil(n/3)+2，并有互补高半区；其 Cor. 5.5 给统一 k\ensuremath{\le}12，Cor. 5.6 给所有奇数 n\ensuremath{\le}27。

\item Hou–Zhao [HZ2]：当 n 不是素数时，Carlet 猜想成立。

\item Carlet–Hou [CH, Table 2]：15,16\ensuremath{\in}K\_31。

\item 互补维数对称：k\ensuremath{\in}\ensuremath{K_{n}} \ensuremath{\Longleftrightarrow} n−k\ensuremath{\in}\ensuremath{K_{n}}。

\end{itemize}

\subsection{奇素数低半区 bookkeeping}

只考虑奇素数 n 和 3\ensuremath{\le}k\ensuremath{\le}(n−1)/2。若 n\ensuremath{\le}27，由 HZ1 已知。n=29 时 ceil(29/3)+2=12，仅剩 k=13,14；附录 A 的两张证书解决。n=31 时阈值为 13，低半区 k\ensuremath{\le}15 中仅剩 14,15；k=14 由附录 A，k=15 由 CH Table 2。对 n\ensuremath{\ge}37，ceil(n/3)+2\ensuremath{\ge}15，故 k\ensuremath{\le}15 全由 HZ1 覆盖；k=16 由命题 9.3；k\ensuremath{\ge}17 由于低半区自动有 n\ensuremath{\ge}2k+1，由命题 9.1、9.2 覆盖。

\begin{quote}\textbf{候选主定理 10.1. 对每个奇数 n 和每个 3\ensuremath{\le}k\ensuremath{\le}n−3，存在 k 维 \ensuremath{F_2}-线性子空间 E\ensuremath{\subset}\ensuremath{\mathbb F_{2^n}} 满足 \ensuremath{\Sigma}\_\{0\ensuremath{\ne}x\ensuremath{\in}E\}\ensuremath{x^{-1}}=0。  [完整证明候选]}\end{quote}

\begin{proof}[证明] 

若 n 非素数，使用 [HZ2]。以下设 n 为奇素数。由互补维数对称，只需处理 3\ensuremath{\le}k\ensuremath{\le}(n−1)/2。n\ensuremath{\le}27 由 [HZ1, Cor. 5.6]。n=29 与 n=31 由上一段列出的显式证书、[CH] 和 HZ1 覆盖。若 n\ensuremath{\ge}37，则 k\ensuremath{\le}15 由 [HZ1, Thm. 5.4]；k=16 由命题 9.3；k\ensuremath{\ge}17 时 n\ensuremath{\ge}2k+1，命题 9.1、9.2 给 \ensuremath{X_{k}}(\ensuremath{\mathbb F_{2^n}})\ensuremath{\ne}\ensuremath{\varnothing}，再由命题 3.2 转回零和 k-subspace。故低半区全部成立，互补对称给高半区。

\end{proof}

\section{附录 A. n=29、31 的显式精确证书}

本附录直接给出闭合有限尾部所需的三张证书。有限域元素以 polynomial basis 的 bitmask 编码：\ensuremath{\Sigma} \ensuremath{a_{i}} \ensuremath{\alpha}\textasciicircum\{\}i 编码为整数 \ensuremath{\Sigma} a\_i2\textasciicircum\{\}i，加法为 XOR。数据与既有证书稿一致。

\subsection{A.1 n=29, k=13}

有限域 \ensuremath{\mathbb F_{2^29}}=\ensuremath{F_2}[\ensuremath{\alpha}]/(\ensuremath{\alpha}\textasciicircum\{\}29+\ensuremath{\ensuremath{\alpha}^2}+1)。基：

\begin{center}\begin{minipage}{0.95\linewidth}\small \texttt{0x00000001 0x0236c002 0x17b15004 0x16f41008 0x0aa45010 0x17b08020 0x178f8040 0x152d9080 0x0b024100 0x1ecb9200 0x11e0c400 0x092bd800 0x0561e000}\end{minipage}\end{center}

对应 \ensuremath{L_{E}}(X)=\ensuremath{\Sigma}\_\{i=0\}\textasciicircum\{\}\{13\}c\_iX\textasciicircum\{\}\{\ensuremath{2^{i}}\} 的系数 \ensuremath{c_0},\ensuremath{\ldots},\ensuremath{c_{13}}：

\begin{center}\begin{minipage}{0.95\linewidth}\small \texttt{0x1a8d2e2f 0x00000000 0x0b39aa9c 0x0f919f54 0x06a764cc 0x0d93d023 0x0cfe406e 0x02d45aa5 0x03f76408 0x1bc649a2 0x10330d40 0x191ffaf2 0x0a266fda 0x00000001}\end{minipage}\end{center}

binary rank = 13；\ensuremath{c_1}=0；8191 个非零元素的直接 reciprocal sum = 0。

\subsection{A.2 n=29, k=14}

同一有限域中，基：

\begin{center}\begin{minipage}{0.95\linewidth}\small \texttt{0x00000001 0x00538402 0x1abd8404 0x18eb0408 0x1aec0010 0x0d550020 0x15c60040 0x1f908480 0x14320500 0x04958200 0x13470800 0x0f4a1000 0x0f172000 0x14fe4000}\end{minipage}\end{center}

系数 \ensuremath{d_0},\ensuremath{\ldots},\ensuremath{d_{14}}：

\begin{center}\begin{minipage}{0.95\linewidth}\small \texttt{0x19ce08e7 0x00000000 0x08d8433a 0x04bd1933 0x121c9c05 0x1b5ac54f 0x046f1aa3 0x1b021f2d 0x1674a568 0x0eee40d5 0x037b879d 0x010d3c06 0x1b2e2c21 0x02427c2c 0x00000001}\end{minipage}\end{center}

binary rank = 14；\ensuremath{d_1}=0；16383 个非零元素的直接 reciprocal sum = 0。

\subsection{A.3 n=31, k=14}

有限域 \ensuremath{\mathbb F_{2^31}}=\ensuremath{F_2}[\ensuremath{\beta}]/(\ensuremath{\beta}\textasciicircum\{\}31+\ensuremath{\ensuremath{\beta}^3}+1)。基：

\begin{center}\begin{minipage}{0.95\linewidth}\small \texttt{0x00000001 0x3b088937 0x171012ec 0x6afce627 0x5e288246 0x4d96a9a8 0x4bf132d0 0x59374007 0x39117efd 0x460fc715 0x1beb9677 0x566e2dcd 0x19437aab 0x00058364}\end{minipage}\end{center}

系数：

\begin{center}\begin{minipage}{0.95\linewidth}\small \texttt{0x1b99e28a 0x00000000 0x435fe9c9 0x04ecca32 0x148e3fe3 0x79802dca 0x32fcb2ff 0x0019ea0d 0x3038b402 0x7f441af6 0x18ecd503 0x347b8d50 0x285c194d 0x48766441 0x00000001}\end{minipage}\end{center}

binary rank = 14；\ensuremath{X^2} coefficient = 0；16383 个非零元素的直接 reciprocal sum = 0。

\section{附录 B. 验证算法与复现实验}

\subsection{B.1 Subspace polynomial verifier}

程序从 \ensuremath{L_{0}}(X)=X 开始，对 basis vector v 逐个应用 L\ensuremath{^+}(X)=L(X)²+L(v)L(X)。由于所有系数位于二元有限域，平方通过有限域乘法完成。最终检查完整 coefficient list 与附录 A 一致，并确认第二个 coefficient（\ensuremath{X^2} 项）为 0。

\subsection{B.2 Exhaustive reciprocal sum}

程序使用 Gray code 枚举 span 中全部非零元素，相邻元素只翻转一个 basis vector。对每个非零 x，用 \ensuremath{x^{2^n−2}} 计算逆元并 XOR 累加。该路线在逻辑上不依赖 Lemma 2.2 的 coefficient criterion，因此是独立验证。

\subsection{B.3 Rabin irreducibility}

29、31 均为素数。对 degree-n binary polynomial p，Rabin criterion 在此只需验证 \ensuremath{X^{2^n}}\ensuremath{\equiv}X (mod p) 以及 gcd(p,\ensuremath{X^2}−X)=1。附带程序逐次平方精确检查。

Supplementary file: verify\_carlet\_paper.py。

\section{附录 C. 数值门槛与独立整数检查}

\subsection{C.1 参数表}

\begin{center}\small\begin{tabular}{|l|l|l|l|}\hline
k & \ensuremath{N_{k}} & lo\ensuremath{g_2} C(\ensuremath{N_{k}}−1,k−2)（约） & 用途 \\\hline16 & 165 & 65.8390 & n\ensuremath{\ge}37 单独处理 \\\hline17 & 184 & 71.6348 & n\ensuremath{\ge}35 gateway \\\hline18 & 204 & 77.5198 & 解析区间起点 \\\hline\end{tabular}\end{center}

\subsection{C.2 精确门槛}

\begin{itemize}

\item C(183,15)<\ensuremath{2^{72}}，因此 \ensuremath{\beta}\_17<\ensuremath{2^{258}}；在 n=35 时 middle ratio<\ensuremath{2^{-4.5}}<1/16。

\item C(164,14)<\ensuremath{2^{66}}，因此 \ensuremath{\beta}\_16<\ensuremath{2^{233}}；在 n=37 时 middle ratio<\ensuremath{2^{-26}}。

\item k\ensuremath{\ge}18 时解析指数 −\ensuremath{k^2}/4+(9/2)k−2 在 k=18 等于 −2，并严格递减。

\end{itemize}

附带程序用整数平方比较验证含半整数指数的 k=17 门槛，避免任何浮点误差。

\section{参考文献}

[Carlet1] C. Carlet, Two Generalizations of Almost Perfect Nonlinearity, Journal of Cryptology 38 (2025), Article 20. DOI 10.1007/s00145-025-09538-5.

[Carlet2] C. Carlet, On the vector subspaces of \ensuremath{\mathbb F_{2^n}} over which the multiplicative inverse function sums to zero, Designs, Codes and Cryptography 93 (2025), 1237–1254. DOI 10.1007/s10623-024-01531-6.

[CH] C. Carlet and X.-D. Hou, More on the sum-freedom of the multiplicative inverse function, Designs, Codes and Cryptography 93 (2025), 4981–4997. DOI 10.1007/s10623-025-01712-x; arXiv:2407.14660.

[HZ1] X.-D. Hou and S. Zhao, Two absolutely irreducible polynomials over \ensuremath{F_2} and their applications to a conjecture by Carlet, Finite Fields and Their Applications 109 (2026), 102713. DOI 10.1016/j.ffa.2025.102713; arXiv:2502.04545.

[HZ2] X.-D. Hou and S. Zhao, On a conjecture about the sum-freedom of the binary multiplicative inverse function, Designs, Codes and Cryptography 94 (2026), Article 86. DOI 10.1007/s10623-026-01819-9; arXiv:2504.21805.

[GT] Z. Gemmell and T. Trudgian, Applying hypersurface bounds to a conjecture by Carlet, arXiv:2512.05431 (2025).

[PS] R. Pink and S. Schieder, Compactification of a Drinfeld Period Domain over a Finite Field, arXiv:1007.4796.

[EPW] B. Elias, N. Proudfoot and M. Wakefield, The Kazhdan–Lusztig polynomial of a matroid, Advances in Mathematics 299 (2016), 36–70; arXiv:1412.7408.

[WZ] D. Wan and D. Zhang, Betti number bounds for varieties and exponential sums, Advances in Mathematics 490 (2026), 110852. DOI 10.1016/j.aim.2026.110852; arXiv:2501.12623. 本文 weak Lefschetz 引用 arXiv 版本 Lemma 2.3.1，Betti bound 引用 Corollary 4.2.4.

[Stacks] The Stacks Project, Section 59.45, Topological invariance of the small étale site; Proposition 59.45.4, Tag 03SI.

[BBD] A. Beilinson, J. Bernstein and P. Deligne, Faisceaux pervers, Astérisque 100 (1982).

\section{后续研究说明}

本文底稿形成于 2026 年 8 月。2026 年 10 月 5 日，Kaimin Cheng 在 arXiv 发布预印本 \emph{Counting zero-sum subspaces for the multiplicative inverse function}，摘要声称结合计数递推与 Fourier 分析证明 Carlet 猜想。本文没有审计该预印本的证明，也不据此作优先权判断。该公开结果应纳入后续文献讨论。本文的几何论证仍按候选证明呈现，须经独立逐行审查。

\section*{参考文献补充}

[Cheng] K. Cheng, \emph{Counting zero-sum subspaces for the multiplicative inverse function}, arXiv:2610.06445v1, submitted 5 October 2026. https://arxiv.org/abs/2610.06445.\end{document}